\documentclass{article}
\usepackage[utf8]{inputenc}
\usepackage[T1]{fontenc}
\usepackage{booktabs}
\usepackage{siunitx}
\usepackage{xintexpr}
\usepackage{pgffor} % Para los bucles \foreach

% Configuramos siunitx
\sisetup{
	round-mode = places,
	round-precision = 8, % <--- Cambia este número según desees
	scientific-notation = fixed,
	fixed-exponent = 0,
	table-format = 1.8 % Ayuda a alinear las columnas S
}

% Definición robusta con xintfloatexpr
% #1 = x, #2 = n
\newcommand{\trxn}[2]{%
	\xintthefloatexpr
	% Definimos las variables para claridad
	% dx = #1/#2
	% g(t) = exp(-(t^2))
	((#1)/(#2)) * (
	(exp(-(0)^2) + exp(-(#1)^2))/2 + 
	add(exp(-(i*(#1)/(#2))^2), i=1..#2-1)
	)
	\relax
}

\begin{document}	
	\begin{table}[h]
		\centering
		\caption{Valores aproximados de $f(x) = \int_0^x e^{-t^2} dt$} 
		\vspace{1em}
		\begin{tabular}{c *{5}{S}}
			\toprule
			& \multicolumn{5}{c}{Iteraciones ($n$)} \\
			\cmidrule(lr){2-6}
			$x$ & {8} & {10} & {14} & {16} & {20} \\
			\midrule
			0.0 & {\trxn{0}{8}}  & {\trxn{0}{10}}  & {\trxn{0}{14}}  & {\trxn{0}{16}}  & {\trxn{0}{20}} \\
			0.5 & {\trxn{0.5}{8}}& {\trxn{0.5}{10}}& {\trxn{0.5}{14}}& {\trxn{0.5}{16}}& {\trxn{0.5}{20}} \\
			1.0 & {\trxn{1.0}{8}}& {\trxn{1.0}{10}}& {\trxn{1.0}{14}}& {\trxn{1.0}{16}}& {\trxn{1.0}{20}} \\
			1.5 & {\trxn{1.5}{8}}& {\trxn{1.5}{10}}& {\trxn{1.5}{14}}& {\trxn{1.5}{16}}& {\trxn{1.5}{20}} \\
			2.0 & {\trxn{2.0}{8}}& {\trxn{2.0}{10}}& {\trxn{2.0}{14}}& {\trxn{2.0}{16}}& {\trxn{2.0}{20}} \\
			\bottomrule
		\end{tabular}
	\end{table}	
	
	
	
\section*{Tabla Automatizada con \texttt{pgffor}}
	
	% 1. Definimos una macro para ir guardando las filas
	\newcommand{\cuerpotabla}{}
	
	% 2. Construimos las filas con un bucle anidado
	\foreach \i in {0, 0.5, 1.0, 1.5, 2.0} {
		% Guardamos el valor de x en la primera columna
		\xdef\filatemp{\i} 
		
		% Bucle para las columnas (n = 8, 10, 14, 16, 20)
		\foreach \n in {8, 10, 14, 16, 20} {
			% Calculamos el valor y lo añadimos a la fila con un &
			\xdef\valor{\trxn{\i}{\n}}
			\xdef\filatemp{\filatemp & \valor}
		}
		
		% Añadimos la fila completa al cuerpo de la tabla y cerramos con \\
		\xdef\cuerpotabla{\cuerpotabla \filatemp \\}
	}
	
	\begin{table}[h]
		\centering
		\caption{Resultados automatizados para $n=\{8, 10, 14, 16, 20\}$}
		\vspace{1em}
		\begin{tabular}{c *{5}{S}}
			\toprule
			$x$ & {$n=8$} & {$n=10$} & {$n=14$} & {$n=16$} & {$n=20$} \\
			\midrule
			\cuerpotabla
			\bottomrule
		\end{tabular}
	\end{table}
\end{document}